\section{See What Humans See：先理解人类视觉为何如此重要}

\subsection{视觉是生物智能最重要的入口之一}

讲者从进化史开始讲视觉，不是为了增加背景故事，而是为了回答一个很根本的问题：\textbf{为什么计算机视觉会成为 AI 最早、也最核心的子方向之一。} 早在约 5.4 亿年前，动物界出现最初的感光细胞之后，视觉就不再只是一个``传感器''，而是推动生物体快速识别环境、躲避危险、寻找食物、识别同伴的关键能力。\footnote{根据字幕 00:01:15--00:03:05，讲者以 Cambrian explosion 与 Andrew Parker 的视觉军备竞赛观点解释视觉的进化意义。}

这段回顾对应着一个研究启示：\textbf{视觉不是普通输入模态，它天然携带了大量压缩后的世界结构。} 人类依赖视觉完成生存、工作、娱乐、社交与学习，这意味着如果机器也能稳定地解析视觉世界，它就有机会触碰到更一般的智能。

\begin{figure}[H]
\centering
\begin{tikzpicture}[every node/.style={font=\small, align=center}]
\node[draw, rounded corners, fill=blue!10, minimum width=2.7cm, minimum height=1cm] (evo) at (0,0) {光感受器出现\\环境感知开始};
\node[draw, rounded corners, fill=blue!16, minimum width=2.9cm, minimum height=1cm] (arms) at (4.2,0) {视觉军备竞赛\\识别能力快速演化};
\node[draw, rounded corners, fill=blue!22, minimum width=3.0cm, minimum height=1cm] (human) at (8.6,0) {人类视觉成为\\主要感知与认知入口};
\node[draw, rounded corners, fill=blue!8, minimum width=3.4cm, minimum height=1cm] (cv) at (4.3,-2.0) {计算机视觉尝试把这种能力\\转化为机器可学习的问题};
\draw[-{Latex[length=3mm]}, thick] (evo) -- (arms);
\draw[-{Latex[length=3mm]}, thick] (arms) -- (human);
\draw[-{Latex[length=3mm]}, thick] (human) -- (cv);
\end{tikzpicture}
\caption{从进化视角理解视觉智能：视觉先是生存能力，随后才成为机器学习问题\protect\footnotemark}
\end{figure}
\footnotetext{根据字幕 00:01:15--00:03:05 整理。讲者强调视觉既是进化史中的关键转折，也是现代 AI 的起点之一。}

\subsection{对象识别为什么看似容易、实则极难}

人类看一张图片，往往几乎不费力就能说出``这是一只狗''、``那是一辆车''。但讲者紧接着提醒：\textbf{这种 effortless 感受会误导我们低估对象识别的计算难度。} 从数学上看，同一对象会随着光照、纹理、遮挡、背景、视角、尺度而发生近乎无限的外观变化；从建模上看，模型既要对这些变化保持鲁棒，又要在类别边界上足够敏感。\footnote{根据字幕 00:07:13--00:08:15，讲者用 lighting、texture、background、occlusion、viewing angle、scaling 等变化解释 object recognition 的根本困难。}

\begin{knowledgebox}{对象识别的真正难点}
\begin{itemize}
    \item \textbf{类内变化巨大}：同一类别在现实世界里几乎没有标准外观。
    \item \textbf{类间差异微妙}：很多错误并不是把猫认成车，而是把近邻类别混掉。
    \item \textbf{环境扰动复杂}：背景、遮挡、模糊、阴影都会破坏模式稳定性。
    \item \textbf{监督信号稀疏}：标注只给类别名，却不告诉模型什么特征才该保留。
\end{itemize}
\end{knowledgebox}

这也解释了为什么整个计算机视觉领域会长期把对象识别当作``基础构件''。如果一台机器连图中最基本的对象都无法稳定识别，就更谈不上关系理解、场景推理、视频分析或具身行动。

\subsection{前深度学习时代：从心理学启发到统计学习}

讲者把 pre-deep-learning 的对象识别史分成两大波。第一波非常``人类中心''：研究者会先内省自己如何看世界，然后猜测视觉系统也许是在识别若干部件、形状或几何原语，再把它们组合成对象。\footnote{根据字幕 00:08:18--00:09:23，讲者回顾 70s--90s 基于 parts / shapes composition 的一系列 heroic attempts。}

这一思路的优点是解释性强，缺点是对真实世界的变化过于脆弱。它能处理``规范对象''，但很难覆盖自然场景里的外观复杂性。因此这类模型常常在理论上优雅，在开放环境里失败。

第二波是统计机器学习的进入。它的重要性不只在于用了 SVM、Bayesian 模型或 random field，而在于研究范式变化了：

\begin{enumerate}
    \item 不再先写死识别规则，而是让模型从数据中估计参数。
    \item 不再只追求单个例子的正确解释，而是追求总体分布上的泛化。
    \item 不再把视觉问题当成几何拼图，而是把它当成学习问题。
\end{enumerate}

\begin{table}[H]
\centering
\begin{tabular}{p{3cm}p{4cm}p{5.5cm}}
\toprule
\textbf{阶段} & \textbf{核心假设} & \textbf{主要瓶颈} \\
\midrule
部件/形状组合 & 对象可以由有限几何部件拼装出来 & 面对真实世界外观变化时脆弱，覆盖面不足 \\
统计学习阶段 & 用参数化模型从数据中学习判别边界 & 数据规模有限，特征表达能力不足 \\
深度学习阶段 & 表示学习与大规模监督共同发挥作用 & 对算力、数据、标注体系提出更高要求 \\
\bottomrule
\end{tabular}
\caption{对象识别研究范式的三次跃迁：从规则解释到统计学习，再到大规模表示学习}
\end{table}

\begin{warningbox}{为什么第一波方法``看起来对''却跑不通}
因为人类解释自己认知过程时，经常只看到了结果，没有看见底层表征如何形成。``对象由部件组成'' 在概念上没错，但这句话距离一个可扩展、可鲁棒的识别系统还差了很多层。
\end{warningbox}

\subsection{ImageNet 改变了问题规模，也改变了研究节奏}

计算机视觉真正进入现代阶段，靠的不是某一个巧妙 trick，而是\textbf{数据、算力和模型在同一时间到位了。} 讲者把 ImageNet 放在极其关键的位置：它把``机器能认多少类''这个问题从小规模实验室玩具，推进到接近真实世界复杂度的量级。

ImageNet 的意义至少有三层：

\begin{itemize}
    \item 它提供了前所未有的大规模视觉词汇表。
    \item 它提供了公开 benchmark，使算法比较进入同一坐标系。
    \item 它让深网络终于有足够多的数据去学到层级表示。
\end{itemize}

2012 年之后，AlexNet 成为拐点，不是因为卷积本身突然诞生，而是因为卷积网络、GPU 训练与大规模监督终于形成共振。之后整个视觉领域迅速从``手工特征 + 传统分类器''转向``端到端表示学习''。

\begin{importantbox}{2012 年真正改变了什么}
不是简单地把分类精度提高了几个点，而是改变了研究者对``视觉能不能靠学习直接解决''的信念。自此以后，识别、检测、分割、描述、生成与多模态模型都开始受益于同一种大规模学习范式。
\end{importantbox}

\subsection{视觉理解不会停在对象标签}

这节课最有价值的一点，是它没有把``识别对象''当成终点。讲者明确往后推了好几层：\textbf{对象只是开始，关系、属性、语言、时间和行动才是更完整的视觉智能。}

为什么这么说？因为单个对象标签往往不足以描述真实场景。``图里有一个人、一只羊驼''和``人骑着羊驼''是完全不同的语义层次。真正能支撑搜索、描述、问答和机器人决策的，不是孤立名词，而是结构化场景理解。

\begin{figure}[H]
\centering
\begin{tikzpicture}[every node/.style={font=\small, align=center}, node distance=1.2cm]
\node[draw, rounded corners, fill=blue!8, minimum width=2.2cm, minimum height=0.9cm] (obj) {对象\\Object};
\node[draw, rounded corners, fill=blue!12, right=of obj, minimum width=2.2cm, minimum height=0.9cm] (attr) {属性\\Attribute};
\node[draw, rounded corners, fill=blue!16, right=of attr, minimum width=2.2cm, minimum height=0.9cm] (rel) {关系\\Relation};
\node[draw, rounded corners, fill=blue!20, right=of rel, minimum width=2.2cm, minimum height=0.9cm] (lang) {语言描述\\Caption / QA};
\node[draw, rounded corners, fill=blue!24, right=of lang, minimum width=2.2cm, minimum height=0.9cm] (action) {行动决策\\Action};
\draw[-{Latex[length=3mm]}, thick] (obj) -- (attr);
\draw[-{Latex[length=3mm]}, thick] (attr) -- (rel);
\draw[-{Latex[length=3mm]}, thick] (rel) -- (lang);
\draw[-{Latex[length=3mm]}, thick] (lang) -- (action);
\end{tikzpicture}
\caption{视觉任务的语义升级路径：对象识别只是第一层，更高层次是关系、语言与行动}
\end{figure}

\subsection{从 scene graph 到 captioning，再到视频理解}

讲者回顾了这一系列扩展任务的内在逻辑：

\begin{enumerate}
    \item \textbf{Scene graph / relationship understanding}：让模型不只列对象名词，而是刻画对象之间的结构关系。
    \item \textbf{Image captioning / storytelling}：让视觉表示进入语言空间，支持更自然的人机接口。
    \item \textbf{Dynamic scene understanding}：让模型处理对象运动、交互变化、时序依赖与未来预测。
\end{enumerate}

这一系列任务难度层层上升。对象识别只需要回答``有什么''；关系理解要回答``彼此如何关联''；captioning 要回答``如何用自然语言压缩这个场景''；视频理解则要再进一步回答``谁在什么时候做了什么，以及接下来可能发生什么''。

\begin{knowledgebox}{为什么视频理解比图像理解更难}
\begin{itemize}
    \item 它需要同时保持空间结构和时间结构。
    \item 它必须区分短时动作与长时情节。
    \item 它要求模型处理主体交互，而不仅是主体存在。
    \item 它离真实世界更近，因此也更接近后续的机器人与 agent 问题。
\end{itemize}
\end{knowledgebox}

\subsection{Compositional generalization：为什么 scene graph 很关键}

字幕里有一段很值得单独拆开看。讲者引用心理学家的观点说，\textbf{复杂自然场景中的对象关系，必须作为理解的一部分被编码。} 这句话看似简单，实际上是在指出分类时代的局限：只知道 ``person'' 和 ``horse'' 同时出现，还远远不足以恢复图像语义；真正有意义的是 ``person riding horse''、``horse wearing hat'' 这样的结构表达。\footnote{根据字幕 00:15:43--00:18:39，讲者用 Jeremy Wolfe 的观点引出 relationship understanding、scene graph 与 zero-shot unusual relationship。}

scene graph 的好处在于，它把视觉理解拆成了更可组合的单元：

\begin{itemize}
    \item 对象节点：谁在场景中出现。
    \item 属性节点：对象有什么状态或特征。
    \item 关系边：对象之间以何种方式相连。
\end{itemize}

这种表示直接带来一种非常重要的能力：\textbf{组合泛化。} 如果模型见过 ``person wearing hat'' 和 ``horse standing''，那么它未必需要大量见过 ``horse wearing hat'' 的训练样本，也有机会通过组合规则去理解这个不常见关系。讲者举的例子很典型：``person sitting on chair'' 和 ``fire hydrant on the lawn'' 常见，但 ``person sitting on fire hydrant'' 就很少见；scene graph 之类的组合表示，正是为了让模型在长尾关系上也能做出合理推断。

\begin{figure}[H]
\centering
\begin{tikzpicture}[every node/.style={font=\small, align=center}]
\node[draw, rounded corners, fill=blue!10, minimum width=2.4cm, minimum height=0.9cm] (obj1) at (0,0) {person};
\node[draw, rounded corners, fill=blue!10, minimum width=2.4cm, minimum height=0.9cm] (obj2) at (4.0,0) {horse};
\node[draw, rounded corners, fill=blue!10, minimum width=2.4cm, minimum height=0.9cm] (obj3) at (8.0,0) {hat};
\node[draw, rounded corners, fill=orange!18, minimum width=2.2cm, minimum height=0.8cm] (r1) at (2.0,-1.8) {riding};
\node[draw, rounded corners, fill=orange!18, minimum width=2.2cm, minimum height=0.8cm] (r2) at (6.0,-1.8) {wearing};
\draw[-{Latex[length=3mm]}, thick] (obj1) -- (r1);
\draw[-{Latex[length=3mm]}, thick] (r1) -- (obj2);
\draw[-{Latex[length=3mm]}, thick] (obj2) -- (r2);
\draw[-{Latex[length=3mm]}, thick] (r2) -- (obj3);
\end{tikzpicture}
\caption{scene graph 的核心思想：把名词、属性与关系显式结构化，支持长尾组合推断}
\end{figure}

\begin{importantbox}{这一步为什么重要}
因为很多真实世界错误都不是 ``看不见对象''，而是 ``看见了对象，但没理解它们之间的关系''。从安全场景到具身智能，关系理解都比孤立分类更接近决策所需的语义。
\end{importantbox}

\subsection{Captioning 与 storytelling：视觉第一次真正接上语言}

讲者随后回顾了 2014--2018 这一波工作：用 CNN 提取图像表示，再接上 LSTM 之类的语言模型，实现 image captioning、dense captioning 与 storytelling。\footnote{根据字幕 00:18:39--00:20:23，讲者提到 Andrej Karpathy 等人在 CNN + LSTM 时代推动 image captioning、dense captioning 和视觉叙事任务。}

这件事在今天的多模态 LLM 时代看起来已经很自然，但在当时意义非常大。因为它第一次大规模证明了：\textbf{视觉系统不必只输出标签，也可以输出一句话、一段描述，甚至一组区域级说明。} 这等于把视觉模型从 ``分类器'' 推进到了 ``语言接口''。

\begin{table}[H]
\centering
\begin{tabular}{p{3.2cm}p{4.4cm}p{6.0cm}}
\toprule
\textbf{任务} & \textbf{输出形式} & \textbf{能力提升} \\
\midrule
图像分类 & 单个类别标签 & 回答 ``这是什么'' \\
关系理解 & 结构化对象关系 & 回答 ``谁和谁如何关联'' \\
图像描述 & 一句或多句自然语言 & 回答 ``这幅图在发生什么'' \\
Dense captioning & 区域级语言描述 & 回答 ``图中各局部各自在做什么'' \\
视频叙事 & 时序化语言输出 & 回答 ``一段时间内事件如何展开'' \\
\bottomrule
\end{tabular}
\caption{从标签到语言的升级：视觉系统的输出空间越来越接近人类交流方式}
\end{table}

这一步还带来一个长期影响：一旦视觉表示进入语言空间，它就天然可以和问答、检索、对话、规划与 agent 系统连起来。也正因为如此，captioning 并不只是一个中间时代任务，而是后来多模态大模型的直接前身之一。

\subsection{这一部分真正想说明什么}

``See what humans see'' 并不是一句空话，它代表视觉 AI 的第一阶段目标：把人类已经擅长的视觉感知任务，系统地变成机器可学习、可扩展、可比较的问题。其核心结论可以概括成三点：

\begin{itemize}
    \item 视觉之所以重要，是因为它承载了大量关于世界结构的先验。
    \item 对象识别之所以难，是因为真实世界变化无限，而监督信号有限。
    \item 现代视觉突破不是单点发明，而是数据、算力、表示学习范式共同成熟的结果。
\end{itemize}

\begin{importantbox}{第一部分的收束}
如果说课程前半程回答的是``机器怎样看见东西''，那么这一部分回答的是``为什么视觉本身就值得成为 AI 的主战场''。这为后面讨论 superhuman perception、人类盲点和 human-centered AI 提供了历史与方法论基础。
\end{importantbox}

